\section{Day 9: Continuous Random Variables (Oct 01, 2025)}

Continuous random variables have the property that $P(X \leq a) = P(x < a)$, since by definition $P(X = x)$. We define 3 new distributions

\begin{itemize}
\item For $L < R \in \mathbb{R}$, we can define the Uniform[$L$, $R$] distribution with the density function
\[
f_X(x) = \begin{cases}
    0 & x < L \\
    \frac{1}{R-L} & L \leq x \leq R \\
    0 & x > R
\end{cases}
\]
Notice that whenever $L \leq a \leq b \leq R$, $P(a \leq X \leq b) = \frac{b-a}{R - L}$. Also $P(L \leq X \leq R) = 1$, so we say that this distribution is \textbf{bounded}.
    \item Exponential($\lambda$), for some $\lambda > 0$, where
    \[
    f(x) = \begin{cases}
        \lambda e^{-\lambda x} & x \geq 0 \\
        0 & x < 0 
    \end{cases}
    \]
\end{itemize}

\begin{simplethm}
    Exponential($\lambda$) has the memoryless property.
\end{simplethm}

\noindent Consider the function
\[
\phi(x) = \frac{1}{\sqrt{2 \pi}} e^{- \frac{x^2}{2}}
\]
$\phi(x) \geq 0$, but it is less clear than that $\int_ \infty ^ \infty \phi(x) \, dx = 1$, which is a result proven in the textbook. For now, assume it a density. If $X$ has density $\phi$, we say that $X$ has the Normal(0, 1) or N(0, 1) distribution. Then for all $a \leq b$,
\[
P(a \leq X \leq b) = \int_a^b \phi(x) \, dx = \int_a^b \frac{1}{\sqrt{2\pi}} e^{- \frac{x^2}{2}} \, dx
\]
We claim that we cannot compute this integral analytically, so you would approximate this using software, or tables.

\begin{itemize}
\item Define the Normal($\mu$, $\sigma^2$) distribution, where for any $\mu \in \mathbb{R}$ and $\sigma > 0$, have
\[
f(x) = \frac{1}{\sigma \sqrt{2\pi}} e^{- \frac{(x-\mu)^2}{2 \sigma^2}} = \frac{1}{\sigma} \phi(\frac{x - \mu}{\sigma})
\]
\end{itemize}

\noindent The curve is centered at $\mu$. A larger $\sigma$ means a wider, shorter and `fatter' curve, while a smaller $\sigma$ means a narrower, taller, and `thinner' curve. This is the key distribution for the Central Limit Theorem, and arises naturally when there are `lots of small influences'. \\

\noindent \textbf{Midterm Info}: The midterm covers everything up to and including the discrete distributions, meaning that there are no continuous distributions on the midterm.
